% Auto-generated by e20_failure.py — NeuroPLC Failure Analysis
\section{When Does NeuroPLC Struggle?}
\label{sec:failure}

We systematically characterize the conditions under which NeuroPLC's
three correctness mechanisms degrade, establishing clear applicability
boundaries.

\subsection{Doubleton Arithmetic Degradation}

DA provides its strongest advantage when weight signs are imbalanced across network edges. Over 10 random KAN architectures, the DA/IA ratio ranges from 1.1$\times$ to 1.7$\times$ (mean 1.4$\times$). The ratio correlates with $|\text{sign bias}|$ (Pearson $r=-0.99$). In the worst case (balanced weight signs), DA degrades to IA (ratio $\rightarrow$ 1.0). However, trained KANs naturally develop sign structure through gradient-based optimization: the CWRU-trained KAN [28,16,4] achieves a 3.1$\times$ DA advantage, significantly above the random-weight baseline.

\subsection{Adaptive LUT Degradation}

Adaptive (curvature-aware) LUT sampling provides measurable improvement only when B-spline basis functions exhibit diverse curvature. Over 5 curvature-diversity trials, adaptive sampling improves accuracy by 0.9$\times$ on average (max 1.1$\times$). When all basis functions have similar curvature ($|\phi''|$ variance $<$ 0.1), adaptive degrades to uniform (100\% of trials show $<$1.2$\times$ improvement). For trained KANs on CWRU, curvature diversity is sufficient to make adaptive LUT beneficial.

\subsection{Depth Lipschitz Explosion}

The network Lipschitz constant $L_{\text{net}}$ grows exponentially with KAN depth. At $L=2$ layers, $L_{\text{net}}$ is well within usable range. The exponential growth rate is approximately $e^{2.42\cdot L}$, with $L_{\text{net}}$ doubling every ~0.3 layers. Beyond depth 5--6, the Lipschitz-based error bound becomes too pessimistic for practical deployment certification. This is an inherent limitation of Lipschitz-based analysis, not unique to NeuroPLC. For deep KANs ($L \geq 5$), we recommend statistical (Monte Carlo) validation as a complement to analytical bounds.

\subsection{Domain Shift Boundaries}

The CWRU $\rightarrow$ XJTU-SY domain shift (MMD$_{\text{RBF}}$ = 0.0906, mean Cohen's $d = 0.132$) is classified as \textbf{Severe — near-random performance without adaptation}. 0/28 features show large shift ($d > 0.5$). This explains the E12 zero-shot result and establishes clear deployment preconditions: (1) same sensor type, (2) same sampling rate, (3) similar operating conditions. When these preconditions are violated, fine-tuning on target domain data (or domain adaptation) is required.

\subsection{Summary of Applicability Boundaries}

\begin{table}[h]
\caption{NeuroPLC applicability boundaries and recommended mitigations.}
\label{tab:boundaries}
\centering
\begin{tabular}{@{}p{2.5cm} p{3.5cm} p{3.5cm}@{}}
\toprule
\textbf{Boundary} & \textbf{Condition} & \textbf{Mitigation} \\
\midrule
DA $\rightarrow$ IA & Balanced weight signs (rare in trained KANs) & Fall back to IA; still provides valid bounds \\
\addlinespace
Adaptive $\rightarrow$ Uniform & Low curvature diversity & Uniform LUT is sufficient; no accuracy loss \\
\addlinespace
$L_{\text{net}}$ explosion & Depth $\geq 6$ layers & Monte Carlo validation; limit depth for safety-critical applications \\
\addlinespace
Domain shift & Cross-sensor, cross-sampling-rate & Fine-tune on target domain; apply domain adaptation \\
\bottomrule
\end{tabular}
\end{table}

